\section{Dos técnicas principales para entrenar LLM de matemáticas}

\subsection{Fine-tuning supervisado (SFT)}

\begin{enumerate}
    \item El modelo base se preentrena con datos a escala de Internet
    \item Continúa el preentrenamiento con documentos matemáticos (StackOverflow, arXiv, etc.) $\to$ modelo matemático base
    \item Se hace fine-tuning con datos de preguntas y respuestas \textbf{que incluyen pasos detallados de solución} (puede incluir llamadas a herramientas, como Python)
\end{enumerate}

\begin{warningbox}{Los datos son el mayor cuello de botella}
La parte más costosa del pipeline de SFT es la recolección de datos: requiere anotación humana, limpieza y conversión de formato, y no puede automatizarse por completo.
\end{warningbox}

\subsection{Aprendizaje por refuerzo (RL)}

Cuando el conjunto de datos solo tiene la respuesta final y no los pasos intermedios, se puede entrenar con RL:
\begin{enumerate}
    \item El modelo genera una solución con pasos intermedios
    \item Se compara la respuesta final con la respuesta verdadera: recompensa 1 si es correcta, recompensa 0 si es incorrecta
    \item Se optimiza el modelo con algoritmos como GRPO para maximizar la recompensa
\end{enumerate}

\begin{importantbox}{La clave de RL: la verificabilidad}
RL depende de una señal de verificación confiable. Las respuestas numéricas se pueden comparar directamente, pero ¿cómo se verifica automáticamente la corrección de \textbf{un problema de demostración}? Aquí es exactamente donde reside el valor de los sistemas formales.
\end{importantbox}

